
\chapter{بحث و محدودیت‌ها}
\label{ch:discussion}

این فصل نتایج فصل~\ref{ch:results} را تفسیر می‌کند. چهار مشاهده کلیدی
به‌تفصیل بررسی می‌شوند و سپس محدودیت‌ها و ملاحظات اخلاقی و حقوقی
برشمرده می‌شوند.

\section{فاصله میان کاهش زیان و تولید معتبر}
\label{sec:disc-gap}

مهم‌ترین یافته این پروژه، الگویی است که در هر شش نسل آزمایشی مسیر مستقیم
تکرار شد: زیان تحت تحمیل معلم به‌طور یکنوا کاهش یافت و در همان حال شرط
اجرای خودرگرسیو هیچ‌گاه برآورده نشد. نسخه دوم روشن‌ترین نمایش این الگو
است، زیرا در آن بهترین رفتار زمان اجرا در گام \N{400} رخ می‌دهد، یعنی
جایی که زیان اعتبارسنجی «بالاتر» است، و با کاهش بیشتر زیان تا گام
\N{2000}، رفتار زمان اجرا به صفر می‌رسد (شکل~\ref{fig:v2}).

سه توضیح ممکن برای این مشاهده وجود دارد و شواهد پروژه به تفکیک آن‌ها کمک
می‌کند.

\paragraph{توضیح یک: نبود ظرفیت یادگیری.} ممکن است آداپتر کم‌رتبه
ظرفیت کافی برای نمایش فارسی نداشته باشد. نسخه \en{v6.2} به‌طور خاص برای
ابطال همین فرضیه طراحی شد: آموزش و ارزیابی روی همان \N{32} نمونه. کاهش
\N{32.20} درصدی زیان متنی درون‌نمونه نشان می‌دهد که آداپتر «ظرفیت»
برازش هدف را دارد. بنابراین این توضیح رد می‌شود.

\paragraph{توضیح دو: اریبی مواجهه.} توضیح سازگارتر با شواهد، همان مفهومی
است که در بخش~\ref{sec:exposure-bias} معرفی شد. در آموزش، مدل همیشه
نشانه‌های متنی و صوتی «درست» پیشین را به‌عنوان شرط دریافت می‌کند و
پیش‌بینی یک گام جلوتر مسئله‌ای آسان است، به‌ویژه در معماری «تک‌گویی درونی»
که متن هم‌تراز‌شده با زمان در ورودی حاضر است. اما در تولید خودرگرسیو، مدل
باید هم متن و هم صوت را از صفر بسازد. اگر نشانه‌های نخست کمی از توزیع
مرجع فاصله بگیرند، مدل به ناحیه‌ای می‌رود که در آموزش ندیده است و خطا
به‌صورت تجمعی رشد می‌کند.

مشاهده‌ای که این توضیح را تقویت می‌کند، در جدول~\ref{tab:v45} دیده می‌شود:
با کاهش زیان، شمار ردیف‌هایی که «اصلاً» گفتاری تولید می‌کنند از
\N{9} به \N{4} یا \N{5} کاهش می‌یابد. یعنی مدل با برازش بهتر هدف، در تولید
مستقل به سکوت میل می‌کند. این رفتار با یک مدل که تنها «به‌قدر کافی
یاد نگرفته» سازگار نیست؛ با مدلی سازگار است که آموخته است در حضور متن
ورودی مرجع به‌درستی پیش‌بینی کند و در غیاب آن فرو می‌ریزد.

نسخه \en{v6.2} کوشید همین را جبران کند: حذف تصادفی زمان‌بندی‌شده ورودی
متنی، سهم نشانه‌های مرجع پنهان‌شده را از \N{0.25} به \N{0.75} افزایش می‌داد
تا مدل بیاموزد بر پیش‌بینی خودش تکیه کند. این مداخله، سیگنال یادگیری را
حفظ کرد اما به تولید معتبر منتقل نشد. توضیح احتمالی آن است که با
\N{32} نمونه آموزشی و \N{200} گام، بودجه در دسترس برای بازآموزی یک رفتار
تولیدی پایدار به‌مراتب کمتر از حد لازم است.

\paragraph{توضیح سه: انگلیسی‌محوری مدل پایه.} مدل پایه به‌شدت
انگلیسی‌محور است. در نسل‌های نخست مشاهده شد که خروجی چک‌پوینت‌های
میانی، گفتار «انگلیسی» است و نه فارسی. یعنی آداپتر نتوانسته توزیع
زبان تولیدی را جابه‌جا کند و تنها بر جزئیات آوایی اثر گذاشته است. تلاش
نسخه چهارم برای آموزش گزینشی درج‌گرهای متنی، پاسخ مستقیم به همین فرضیه بود
و آن نیز کافی نبود.

\paragraph{نتیجه‌گیری روشمند.} فارغ از این‌که کدام توضیح سهم بیشتری دارد،
یک نتیجه عملی قطعی است و به نظر ما مهم‌ترین سهم علمی این نوشتار است:

\begin{quote}
در تطبیق کم‌پارامتر یک مدل گفتار‌به‌گفتار به زبانی جدید، زیان اعتبارسنجی
تحت تحمیل معلم معیار کافی و حتی معیار قابل اتکایی برای انتخاب چک‌پوینت
نیست. آزمون تولید خودرگرسیو در زمان اجرا باید بخشی از معیار انتخاب باشد و
نه یک بررسی پس از انتخاب.
\end{quote}

اگر پروژه معیار انتخاب را تنها بر زیان بنا کرده بود، چک‌پوینت گام
\N{8000} نسخه نخست، با کمترین زیان اعتبارسنجی و با بهبود معنادار نسبت به
مدل پایه (\N{1.7272} در برابر \N{3.5670} روی داده کنارگذاشته)، به‌عنوان
«مدل گفتار‌به‌گفتار فارسی» گزارش می‌شد؛ در حالی که همان چک‌پوینت در سرور
رسمی صوت تقریباً ساکت و هیچ متنی تولید نمی‌کرد.

\section{چرا نرخ خطای واژه بسیار بیشتر از نرخ خطای نویسه است}
\label{sec:disc-cer-wer}

در سنجه بازبازشناسی همان \N{9} خروجی پنل مکانیکی با پاسخ‌گوی
\en{Qwen2.5-0.5B}، نرخ خطای نویسه \N{7.14} درصد و نرخ خطای واژه
\N{30.41} درصد است؛ یعنی نرخ خطای واژه بیش از چهار برابر است. این نسبت در
نگاه اول نگران‌کننده به نظر می‌رسد اما تا حد زیادی یک اثر
«تعریفی» در زبان فارسی است.

سه سازوکار در این فاصله سهم دارند. نخست، \مهم{مرز واژه در فارسی}. در
تعریف به‌کاررفته، نیم‌فاصله یک مرز واژه به شمار می‌آید. اگر سرویس بازشناسی
یک واژه مرکب را با فاصله کامل بنویسد در حالی که مرجع آن را با نیم‌فاصله
نوشته است، یک واژه به دو واژه تبدیل می‌شود و این تنها اختلاف، دو ویرایش
واژه‌ای تولید می‌کند، در حالی که در سطح نویسه تنها یک نویسه تفاوت دارد.
دوم، \مهم{یکسان‌نبودن صورت نوشتاری}. برای بسیاری از واژه‌های فارسی چند
صورت نوشتاری پذیرفته وجود دارد. سوم، \مهم{اثر طول}. میانگین طول واژه در
داده این پنل حدود چهار نویسه است؛ در چنین شرایطی هر خطای یک‌نویسه‌ای
به‌طور خودکار به یک خطای کامل واژه‌ای تبدیل می‌شود.

شاهد پشتیبان این تفسیر در داده هست: در هر \N{9} ردیف، شمار نویسه‌های فرضیه و
مرجع تقریباً یکسان است (برای نمونه \N{52} و \N{52} در ردیف نخست) در حالی
که شمار واژه‌ها تفاوت دارد (\N{11} و \N{12}). یعنی محتوای نویسه‌ای تقریباً
حفظ شده و اختلاف عمدتاً در بخش‌بندی است.

با این حال، دو محدودیت این تفسیر را باید صریح گفت. نخست، بخشی از خطا
می‌تواند از تلفظ نادرست موتور گفتاری بیاید و نه از بخش‌بندی؛ تفکیک این دو
سهم بدون گوش‌دادن انسانی ممکن نیست. دوم، سرویس بازشناسی هم در سامانه و هم
در اندازه‌گیری مشارکت دارد؛ خطای مشترک آن‌ها می‌تواند نتیجه را به هر دو
سمت جابه‌جا کند. بنابراین نرخ خطای نویسه \N{7.14} درصد را نمی‌توان به
«گفتار تولیدی برای انسان کاملاً قابل فهم است» ترجمه کرد.

\section{چرا تطبیق پاسخ‌گو زیان را کم کرد ولی کیفیت را بهتر نکرد}
\label{sec:disc-reference}

آزمایش بخش~\ref{sec:res-lora} یک تناقض ظاهری دارد: تطبیق کم‌پارامتر
پاسخ‌گو زیان توسعه را \N{31.18} درصد کاهش داد اما در ارزیابی معنایی از خط
پایه پیشی نگرفت. پاسخ این تناقض در ماهیت «هدف» آموزشی نهفته است.

هدف آموزشی این آداپتر، بازتولید «نوبت گوینده بعدی» در داده مکالمه بود.
اما نوبت گوینده بعدی در یک برنامه گفت‌وگوی طبیعی، پاسخ یک دستیار به یک
پرسش نیست. آن نوبت می‌تواند یک تأیید کوتاه، یک تغییر موضوع، یک شوخی، یک
پرسش متقابل یا ادامه روایت خود گوینده باشد. مدل با آموزش روی این هدف،
«سبک» نوبت‌گیری طبیعی را بهتر می‌آموزد و زیان کاهش می‌یابد، اما
دستور‌پیروی و پاسخ‌دهی مرتبط را نمی‌آموزد و ممکن است آن را حتی تخریب
کند.

شاهد مستقیم و تعیین‌کننده این تفسیر در آزمون نهایی معنایی است: بازوی
«مرجع هم‌ترازشده خودکار نوبت گوینده بعدی» کمترین امتیاز ارتباط را گرفت،
\N{0.950}، که حتی از خط پایه \N{0.5} میلیارد‌پارامتری (\N{1.425}) نیز
پایین‌تر است (شکل~\ref{fig:semantic}). یعنی داور خودکار، همان هدفی را که
آداپتر کوشیده بود تقلید کند، به‌عنوان پاسخ نامرتبط ارزیابی کرده است.

نتیجه عملی این مشاهده برای پژوهش‌های بعدی روشن است: نوبت‌های
هم‌ترازشده خودکار از مکالمه طبیعی، مواد «ارزشمندی» برای یادگیری
زمان‌بندی گفتار، نوبت‌گیری و پدیده‌های تعاملی‌اند، اما بدون یک مرحله
گزینش و بازنویسی، هدف «ضعیفی» برای آموزش پاسخ‌دهی دستور‌پیرو
هستند. همین نکته احتمالاً بخشی از شکست مسیر مستقیم را نیز توضیح می‌دهد،
زیرا مدل مستقیم روی همان جفت‌ها آموزش دیده است.

\section{تعارض ساختاری میان نیازمندی‌ها}
\label{sec:disc-conflict}

تحلیل جدول~\ref{tab:positioning} یک تعارض ساختاری در تعریف پروژه را آشکار
می‌کند که ارزش ثبت دارد. نیازمندی ن۲ قابلیت تمام‌دوطرفه را می‌خواهد و
نیازمندی ن۳ تأخیر حداکثر \N{500} میلی‌ثانیه، در حالی که نیازمندی ن۶
اجرا روی پردازنده گرافیکی با \N{12} تا \N{24} گیگابایت حافظه را الزام
می‌کند.

مدل‌هایی که قابلیت تمام‌دوطرفه «ذاتی» دارند، یعنی جریان کاربر و جریان
خود را هم‌زمان مدل می‌کنند، در زمان انجام این پروژه در بازه حافظه هدف جای
نمی‌گرفتند. مدل پایه مسیر مستقیم این پروژه در آموزش تا
\N{28.640} گیگابایت حافظه مصرف کرد. در سوی دیگر، معماری آبشاری در بازه
حافظه هدف جای می‌گیرد اما تأخیر آن، همان‌طور که در
جدول~\ref{tab:cascade-rows} دیده می‌شود، با میانه \N{5584.7} میلی‌ثانیه
برای تولید کامل پاسخ، از هدف \N{500} میلی‌ثانیه بسیار دور است.

این یک نقص پیاده‌سازی نیست، بلکه یک محدودیت واقعی حوزه در زمان انجام
پروژه است. راه‌های عبور از آن در فصل~\ref{ch:conclusions} پیشنهاد شده
است، به‌ویژه تولید جویباری که در آن نخستین قطعه صوت پیش از تکمیل پاسخ
پخش می‌شود و بنابراین سنجه نخستین صوت از زمان تولید کامل جدا می‌گردد.

\section{ارزش و مرز نتیجه معنایی خودکار}
\label{sec:disc-judge}

نتیجه آزمون نهایی معنایی، مثبت و از هر هفت شرط ازپیش‌تعریف‌شده عبور کرده
است. طراحی این آزمون چند نقطه قوت روشمند دارد: مرحله صلاحیت جدا از آزمون
نهایی، شاخص ردیف‌ها و آستانه‌های قفل‌شده پیش از اجرا، دوسوکورسازی و
تصادفی‌سازی ترتیب بازوها، کدگشایی قطعی داور، و اعتبار کامل هر \N{240}
فراخوانی با نرخ توافق تکرار \N{1.00}.

با این حال، سه محدودیت آن را از یک نتیجه انسانی جدا می‌کند. نخست،
\مهم{هم‌خانواده‌بودن داور و نامزد}. داور و نامزد از یک خانواده مدل‌اند؛
پیشینه ارزیابی خودکار نشان داده است که داورها به خروجی مدل‌های
هم‌خانواده تمایل نشان می‌دهند~\cite{zheng2023mtbench}. دوم،
\مهم{کالیبره‌نبودن سنجه‌نامه}. سنجه‌نامه امتیازدهی با امتیاز انسانی
فارسی‌زبان کالیبره نشده است؛ بنابراین نمی‌دانیم امتیاز \N{3.150} از
\N{4} به چه سطحی از رضایت کاربر واقعی متناظر است. سوم،
\مهم{محدوده معنایی}. این آزمون تنها ارتباط و انسجام را می‌سنجد و درباره
درستی واقعی محتوا، ایمنی پاسخ یا طبیعی‌بودن گفتار چیزی نمی‌گوید.

با این مرزها، تفسیر درست چنین است: زنجیره آبشاری، پاسخ‌هایی تولید می‌کند
که یک داور خودکار آن‌ها را به‌طور نظام‌مند و با فاصله زیاد
(\N{+1.725} در ارتباط) مرتبط‌تر از خط پایه ارزیابی می‌کند. این یک شاهد
مهندسی معتبر برای کارکرد سامانه است و نه یک ادعای کیفیت ادراکی.

\section{محدودیت‌ها}
\label{sec:limitations}

محدودیت‌های این پروژه به‌صورت صریح و فهرست‌شده چنین است.

\begin{enumerate}
\item \مهم{سخت‌افزار هدف.} ارزیابی روی هیچ پردازنده گرافیکی فیزیکی با
      \N{12} تا \N{24} گیگابایت حافظه انجام نشد. سخت‌افزار موجود
      \N{93.09} گیگابایت حافظه داشت و سقف نرم‌افزاری حافظه، آن را به یک
      کارت واجد شرایط تبدیل نمی‌کند.

\item \مهم{تأخیر سرتاسری.} سنجه $T_{\text{first\_audio}}$ با رویداد
      تأییدشده مرورگر اندازه‌گیری نشد و شمار ردیف‌های واجد شرایط رسمی صفر
      است. زمان تولید کامل پاسخ جانشین آن نیست. ادعای تأخیر کمتر از
      \N{500} میلی‌ثانیه با شواهد موجود مجاز نیست.

\item \مهم{مطالعه انسانی.} مطالعه \N{5} تا \N{10} گویشور فارسی با
      دست‌کم دو نفر بالای \N{60} سال انجام نشد. وضعیت ثبت‌شده یک
      شرکت‌کننده و صفر مجموعه امتیاز کامل است. بنابراین هیچ یافته‌ای
      درباره طبیعی‌بودن گفتار خروجی، رضایت کاربری یا تعامل سالمندان
      فارسی‌زبان قابل ادعا نیست.

\item \مهم{کیفیت معنایی.} کیفیت معنایی تنها با یک تقریب خودکار
      هم‌خانواده سنجیده و مثبت شد. طبیعی‌بودن، درستی تلفظ، رضایت کاربر و
      درستی واقعی محتوا با قضاوت انسانی سنجیده نشد.

\item \مهم{برچسب‌های وقفه.} برچسب‌های مرجع ارزیابی آشکارساز از تفکیک
      گوینده خودکار به دست آمده‌اند و بازبینی مستقل انسانی ندارند. دقت
      نقطه‌ای \N{81.06} درصد از \N{80} درصد بالاتر است، اما کران پایین
      فاصله اطمینان بلوکی در سطح نشست \N{74.44} درصد است. عبور رسمی از
      نیازمندی ن۵ اثبات نشده است.

\item \مهم{بازبینی شنیداری داده.} بازبینی شنیداری چهل‌ردیفی داده و
      بازبینی بیست‌و‌چهار‌ردیفی رخدادهای تعاملی، تحت یک سند انصراف زمانی
      مستند، انجام نشد. هیچ ردیفی از مجموعه داده «تأییدشده توسط انسان»
      نیست.

\item \مهم{مسیر مستقیم.} هیچ‌یک از شش نسل آزمایشی، چک‌پوینت واجد شرایط
      استقرار تولید نکرد. نتیجه این مسیر یک نتیجه منفی با سیگنال یادگیری
      مثبت است و نه یک مدل گفتار‌به‌گفتار فارسی موفق.

\item \مهم{مرورگر و میکروفون فیزیکی.} پیوستگی تمام‌دوطرفه در سطح انتقال
      سرور آزموده شد، اما توقف واقعی منبع صوت در یک مرورگر واقعی با
      میکروفون فیزیکی اندازه‌گیری نشد.

\item \مهم{تعمیم مجموعه داده.} مجموعه داده از چهار کانال برنامه گفت‌وگویی
      مشتق شده است و سبک گفتاری آن‌ها نماینده همه گونه‌های گفتار فارسی
      نیست. صدای کانال دستیار نیز با یک صدای واحد بازسازی شده است.

\item \مهم{بازتولید کامل.} تنسورهای بخشی از چک‌پوینت‌های نامزد
      ارتقا‌نیافته به‌عمد حذف شده‌اند و بازتولید دقیق آن‌ها مستلزم اجرای
      مجدد دستور آموزش قفل‌شده است. همه نتایج علمی، پیکربندی‌ها،
      خروجی‌های زمان اجرا و درهم‌سازی‌ها حفظ شده‌اند.
\end{enumerate}

\section{ملاحظات اخلاقی، حریم خصوصی و حقوقی}
\label{sec:ethics}

\subsection{خاستگاه و حقوق داده}

مجموعه داده از یک آرشیو داخلی محتوای گفتاری فارسی مشتق شده است. استفاده
پژوهشی داخلی از این منابع با تأیید استاد راهنما ثبت شده است. این تأیید
«مجوز بازنشر» نیست: صوت خام، زیرنویس خام و مشتقات آن‌ها در مخزن کد
قرار نگرفته‌اند و در بسته تحویل پروژه نیز منتشر نمی‌شوند. عبارت
«داده عمومی در دسترس، مجوز آموزش یا بازنشر آزاد دارد» نادرست است و در
این نوشتار به کار نرفته است.

\subsection{حریم خصوصی در مطالعه انسانی}

سامانه سه حالت نگه‌داری داده دارد: نگه‌داری پرونده صوتی، نگه‌داری تنها
بردار ویژگی کاهش‌یافته، و نگه‌داری تنها سنجه‌ها. حالت پیش‌فرض پیشنهادی
برای مطالعه انسانی، حالت فقط‌ویژگی است تا گفتار قابل بازسازی شرکت‌کننده
ذخیره نشود. آزمون خودکار نیز بررسی می‌کند که در حالت فقط‌سنجه هیچ پرونده
صوتی نوشته نمی‌شود.

در ارزیابی‌های خودکار نیز اصل کمینه‌سازی داده رعایت شده است: مصنوعات
نتیجه، متن رونویسی، متن پاسخ و توضیح داور را به‌صورت متن آشکار ذخیره
نمی‌کنند و تنها درهم‌سازی و مقادیر تجمیعی نگه می‌دارند. تحلیل توصیفی
پس‌نگر بخش~\ref{sec:res-cascade} نیز به‌طور صریح بدون خواندن متن آشکار
اجرا شد.

\subsection{مجوز و مرزهای آن}

کد منبع و مستندات متعلق به این پروژه با مجوز \en{Apache-2.0} منتشر
می‌شوند. این مجوز «شامل» موارد زیر نیست و شرایط مستقل هر یک همچنان
حاکم است: اجزای نرم‌افزاری شخص ثالث، وزن‌های مدل‌های زبانی و صوتی،
مجموعه داده منبع و مشتقات آن، آداپترهای آموزش‌دیده، رسانه محدودشده و
زیرنویس‌ها، و مواد شرکت‌کنندگان مطالعه انسانی. به‌طور خاص، موتور گفتاری
مورد استفاده مجوز متفاوتی دارد و تعهدات آن جداگانه باقی می‌ماند.

\subsection{کاربرد و سوءکاربرد}

سامانه‌ای که گفتار فارسی طبیعی تولید می‌کند، قابلیت سوءکاربرد دارد. در
شکل کنونی، دو عامل این خطر را محدود می‌کند: صدای خروجی، یک صدای
هم‌آمایش‌شده واحد و شناخته‌شده است و شبیه‌سازی صدای افراد پیاده نشده است؛
و پاسخ‌گوی متنی با یک پرامپت قفل‌شده کار می‌کند. با این حال، پروژه هیچ
سنجش ایمنی محتوا انجام نداده است و مدل پاسخ‌گو می‌تواند محتوای نادرست
تولید کند؛ استفاده از این نمونه اولیه در هر زمینه‌ای که پاسخ نادرست
پیامد جدی دارد، بدون یک لایه سنجش ایمنی مجاز نیست.

\section{عبارت‌هایی که این نوشتار از آن‌ها پرهیز کرده است}
\label{sec:forbidden}

برای شفافیت، هشت عبارت زیر با شواهد این پروژه \مهم{قابل ادعا نیستند}. هر
عبارت در گیومه نقل شده است تا روشن باشد که این‌ها ادعاهای این نوشتار
نیستند، بلکه فهرست صریح چیزهایی‌اند که \مهم{گفته نشده‌اند} و نباید از این
پروژه نتیجه گرفته شوند:

\begin{itemize}
\item «مدل \en{Moshika} با موفقیت برای فارسی تنظیم شد.»
\item «تأخیر سامانه کمتر از \N{500} میلی‌ثانیه است.»
\item «دقت رسمی تشخیص وقفه بیش از \N{80} درصد است.»
\item «خروجی گفتاری برای انسان طبیعی و کاملاً قابل فهم است.»
\item «سامانه برای استفاده سالمندان مناسب است.»
\item «پنل نه‌ردیفی، تعمیم سامانه به گفتار فارسی را اثبات می‌کند.»
\item «داده منبع مجوز بازنشر یا آموزش آزاد دارد.»
\item «مجموعه داده توسط انسان برچسب‌گذاری و تأیید شده است.»
\end{itemize}
